%=====================================================================
\section{Conclusion}
\label{sec:conclusion}

Decision orchestration learns annotation procedures, not model weights. Where a guideline is a sequence of tests over enumerable candidates, as for verbal multiword expressions, sixty gold sentences and a few dollars give a procedure that beats direct LLM annotation by 5--18 F and, in English, the best published systems we know of. Where the output is mostly structure, as in \umr, orchestration helps small models and enables routing, but closed decisions alone do not beat a capable generator given examples. Matched baselines, repeated runs and ablations remain necessary.

\section*{Limitations}
We test two schemes in two languages each, and most configurations were run once. Run-to-run noise (up to 0.6 \F for closed mode and 2.3 \F for a 9B model in fragment mode) is of the same order as some of the differences we report; repeated runs and other document orders would show the variance of reflection. The \umr test sets are small (99--157 sentences) and differ between stages of the work; unifying them is planned (\exper{E1}). The Chinese set H had been inspected in earlier analyses, and the batch seed $\Pi_0$ was derived after inspecting the released revision's scores on it, so all comparisons on H are retrospective. PARSEME English comes from edition 1.1, whose systems predate today's pretrained models; edition 1.2 open-track systems remain far ahead in Chinese, and the second PARSEME language we planned (French) was not run. The PARSEME procedures read the lemmas, parts of speech and dependency relations provided with the data, and our unseen-expression counts differ from the published ones. The seeds were written by the authors, so their quality, and the effort it took to write them, affect every result; we have not measured that effort. Several of the models we use, including all three decision-model services and the frontier LLMs, are commercial APIs whose behavior may change; the release must include shareable predictions, exact versions and permitted cached answers before these runs can be called independently reproducible \todo{check that the release includes them}. The routing pilot comes from 20 sentences, routing on H was run at two thresholds, and the ``human'' in routing will be simulated from gold, which overstates real annotators' accuracy. Finally, AnCast and our fine-grained judge depend on an automatic alignment, and the gold itself contains errors. One gold graph has unbalanced brackets and is skipped by the scorer.

\section*{Ethics Statement}
We use the Chinese \umr annotations supplied in the research checkout and the English annotations of \umr 2.0 (CC BY-SA 4.0). \todo{Verify the exact release, license, cleaned-gold provenance, redistribution and commercial-API use permissions.} PARSEME 1.2 Chinese is released under CC BY-NC-SA 4.0 and PARSEME 1.1 English under the licenses in its README. We collected no new data from people and no personal information beyond what the released documents contain. Our goal is to reduce the cost of producing the first annotations for a new scheme, which could reduce demand for annotation work. We see the method as a way to direct annotators' effort to the decisions that need human judgment; procedures that annotate cheaply could also produce training data at scale, and their outputs should be checked before they are used as gold. The experiments added in this revision cost about \$70 in billed API charges, including pilots and discarded runs; earlier experiments used commercial APIs whose cost and price-date audit is still pending \todo{sum the API costs}.

Generative AI assisted author-side simulated review, manuscript editing, and the implementation and running of the batch-reflection and PARSEME experiments, as described; the authors are responsible for all content. This assistance is not formal ARR peer review.
